\documentclass[11pt]{article}
\usepackage[a4paper,margin=25mm]{geometry}
\usepackage{amsmath,amssymb,amsthm,microtype}
\usepackage[hidelinks]{hyperref}
\newcommand{\E}{\mathbb E}\newcommand{\Pp}{\mathbb P}
\newcommand{\CE}{C_{\rm E}}
\newtheorem{theorem}{Theorem}\newtheorem{corollary}[theorem]{Corollary}
\title{One acquired strip and two noise estimates}
\author{}\date{29 October 2026}
\begin{document}\maketitle

The new comparison has one probability interval and one support strip.
Below noise $1/100$, the loss at a single count dominates the entire incoming
tail by fourth moments. Above that point, five neighboring counts retain
enough crossing loss, and four rational endpoint checks finish the argument.
The acquisition reaches this common strip directly from the coarse defect
budget and two contact rays. This note gives the exact joins and the remaining
conditional boundary. No claim about arbitrary laws follows from a strip
theorem without the acquisition step.

For a centered variance-one marginal $P$ and its independent copies, write
\[
\mathcal E(P,n,H)=\sqrt n\{\Pp(X_1+\cdots+X_n\le H)
                         -\Phi(H/\sqrt n)\}-\CE\E|X|^3,
\qquad \CE=\frac{3+\sqrt{10}}{6\sqrt{2\pi}}.
\]
Here $H$ is a raw sum threshold. In particular its normal argument is
$H/\sqrt n$, not $H$.

\begin{theorem}[Complete local comparison]\label{th:local}
Let $n\ge700$ be an integer and $Y=B+U$, where $B$ is Bernoulli of parameter
$p\in[3/8,93/200]$ and $\E(U\mid B)=0$. Suppose the conditional noises obey
\[
 -3/5\le U_0\le a_0,\qquad -a_1\le U_1\le3/5,\qquad
 a_0,a_1\ge0,\quad a_0+a_1\le1/250.
\]
Put $v=p(1-p)$, $s=\E U^2$, $\lambda=ns$, and
$X=(Y-p)/\sqrt{v+s}$. If $0<\lambda\le7/20$ and $|H|\le46/25$, then
$\mathcal E(P,n,H)<-\lambda/100$. At $\lambda=0$ one has
$\mathcal E(P,n,H)<0$.
\end{theorem}
\begin{proof}
For $0<\lambda\le1/100$, apply Theorem~1 of
\texttt{small-common.tex}. Choose an integer $k$ so that
$np+H\sqrt{v+s}=k+u$, with $u\in[-3/10,7/10]$.
The interval has length one, and the proof checks that the anchor is interior.
At $C=\CE\ge0$, that theorem bounds the difference between $\mathcal E$
and the signed binomial excess at $k$ by $-\lambda/100$.
Schulz's strict finite-binomial theorem (Theorem~1, 2016 thesis) bounds the
latter excess by zero for every $0<p<1$ and finite $n$.

For $1/100\le\lambda\le7/20$, use the Compact row of Theorem~1 in
\texttt{unified-comparison.tex}. All probability, support, inward and threshold
hypotheses agree literally. Its floor anchor need not equal the small-noise
anchor: both conclusions concern the same $\mathcal E(P,n,H)$.
The shared noise endpoint is included in both statements; there is no phase
or noise gap. At zero noise, $U=0$ almost surely, so the same strict binomial
theorem applies directly. A limit of positive-noise margins is not used.
\end{proof}

\begin{theorem}[Conditional acquired-witness exclusion]\label{th:acquired}
Suppose $n\ge700$ and $(P,H)$ is a compact, centered variance-one attained
fixed-target witness at $\CE$, with $\mathcal E(P,n,H)\ge0$ and true sharp
predecessor $C_{n-1}\le\CE$. Assume the retained closed-support contact and
all-threshold heat bounds hold for this witness. Suppose the complete coarse
entry theorem supplies, for this same oriented law,
\[
 nD\le32.20665378,\qquad |H|\le1.83605680,
 \quad D=(3+\sqrt{10})\rho-\gamma-3J,
\]
where $\rho=\E|X|^3$, $\gamma=\E X^3$ and $J$ is its actual weighted
support-gap sum. Then such a witness cannot exist.
\end{theorem}
\begin{proof}
The scalar theorem in \texttt{acquisition-direct.tex} first establishes an
actual gap across zero, with $p=\Pp(X>0)\in(3/8,93/200)$.
Write $m_j=\E(X\mid B=j)$ and $\ell=m_1-m_0>0$, where
$B=\mathbf1_{X>0}$. Centering gives $m_0=-p\ell$ and $m_1=(1-p)\ell$.
With $U=(X-m_B)/\ell$ and $\theta=\E(X-m_B)^2$, total variance gives
\[
 1=v\ell^2+\theta,\qquad s=\E U^2=\theta/\ell^2,
 \qquad v+s=\ell^{-2},\qquad X=\frac{B+U-p}{\sqrt{v+s}}.
\]
Thus the normalized marginal in Theorem~\ref{th:local} is exactly the
original law, with $\lambda=n\theta/\ell^2<7/20$ and inward sum $d<1/250$.
It is not an optimizer of a relaxation or a separately chosen product law.
The two contact rays in the second acquisition theorem give the common
amplitude $3/5$. The raw threshold obeys $|H|<46/25$. These are every
hypothesis of Theorem~\ref{th:local}, which yields $\mathcal E<0$ and the
contradiction. Neither the old left/right probability panels nor an additional
noise-localization table enters this implication.
\end{proof}

\section*{The global boundary is unchanged}
For any sum $S$,
\[
 \Pp(-S\le-H)-\Phi(-H/\sqrt n)
 =\Phi(H/\sqrt n)-\Pp(S<H)
 \ge\Phi(H/\sqrt n)-\Pp(S\le H).
\]
The gap between the last two expressions is the atom at $H$. This selects
the positive closed-event orientation; it does not transport the signed
$D$ budget to a reflected law. Attainment and entry must apply afresh to
the chosen oriented witness.

If a genuine base establishes $C_j\le\CE$ through $j=41$, and the retained
attainment, contact and entry implications hold at every least failure, then any
first failure of the sharp bound is in $42,\ldots,699$. Minimality supplies
the predecessor needed in Theorem~\ref{th:acquired}. This does not prove an
unconditional cutoff at $700$: an unresolved earlier failure invalidates
that predecessor argument at later rows. A genuine base through $699$, or
an independent argument closing the remaining range, is still needed.

The coarse Fourier certificate, its two coarse and four signed entry passes,
and their internal cover obligations are retained. The new proof removes
only downstream probability-panel localization, old numerical noise tables,
the separate left/right local union, the Poisson-tail comparisons and the
291-cell negative-inner certificate from this acquired $n\ge700$ route.
The old results at smaller indices and all their proof objects are preserved.
Exact scalar checks and actual-law controls accompany these proofs; they
do not independently establish the retained analytic inputs or the finite base.
\end{document}
